\section{Threats to Validity}
\label{sec:threats}

\paragraph{Construct validity.} Our cost construct $\kappa$ is the
measured value of a cost-like property shipped with each corpus (runtime,
energy, training time). It captures the dominant marginal cost of
acquiring one configuration's measurement but omits fixed costs (build
and deployment machinery, engineer time) and possible parallelism. On
Linux, cost and target coincide (kernel size proxies build effort), which
makes the coupling maximal by construction; we therefore base no
coupling-mechanism claim on Linux and derive the mechanism from the
ENERGY/LCBench contrast, where cost and target are measured
independently. Conversely, the equal-spend budgets ($B = n_{\mathrm{ref}}
\cdot \bar\kappa$) inherit random sampling's expected spend as the
reference; other anchors would shift absolute win rates but not the
ENERGY-versus-LCBench contrast.

\paragraph{Internal validity.} All pipelines are seeded and were verified
to be deterministic end to end (an interrupted Linux sweep reproduced its
completed part bit-identically when rerun). Significance-based results
are FDR-controlled per (corpus, hypothesis) family; the headline H1
counts move only on LCBench (154 $\to$ 144 of 210). Stochastic strategies
use 30--100 seeds; deterministic strategies contribute single runs, so
their per-scope comparisons rest on decision rules rather than tests.
Our \srs{} reimplementation was audited against the artifact
(Section~\ref{sec:srsaudit}); the FSE-artifact corpus itself uses the
authors' shipped samples, insulating the close replication from
reimplementation bias. The artifact's duplicate-containing \srs{} fills
mean its nominal sample sizes are occasionally 1--5 configurations
optimistic; this slightly favors \srs{} in size-indexed comparisons and
is therefore conservative with respect to our cost-indexed conclusions.

\paragraph{External validity.} Linux and LCBench are large measured
pools, not exhaustive configuration spaces; every posterior-risk
statement for them is pool-relative, and pool-construction bias (e.g.,
how TuxKConfig configurations were generated) bounds what the guarantees
mean. LCBench is deliberately outside configurable software; we use it to
probe the framework's domain sensitivity, not to make claims about
hyperparameter optimization practice. Our budget grids emphasize small
samples on the generated families ($n \le 64$, $\le$30\% of pools), the
regime where practitioners with expensive measurements operate but also
the regime least favorable to \srs; the x264 and FSE-artifact corpora
cover the original $t$-wise-scale regime, and the equal-spend gradient
connects the two. We study worst-case (maximum) bounds only; best-case
bounds and the original study's measurement-noise and
restricted-ground-truth robustness experiments are not replicated here.

\paragraph{Conclusion validity.} Counts aggregate scopes, and scopes
within a corpus share systems and pools; we therefore never pool scopes
across corpora for significance testing and report per-corpus counts
throughout. Where we summarize across corpora (e.g., 830/851), the figure
is descriptive, not inferential.
